% !TeX root = ../../geos_mpm_manual.tex
\chapter*{Preface}
\addcontentsline{toc}{chapter}{Preface}
\markboth{Preface}{Preface}
\index{GEOS-MPM}
\index{manual!starter}

This starter manual documents the GEOS Material Point Method implementation as it appears in the uploaded source archive.  It is designed to be regenerated as the code changes: narrative chapters provide the stable structure, while generated appendices collect solver attributes, ParticleFileWriter parameters, particle fields, event attributes, geometry objects, materials, suites, and post-processing tools extracted from the code and schema documentation.

\paragraph{Generation status.} The source archive reports {\ttfamily VERSION\_\allowbreak{}ID = v1.\allowbreak{}1.\allowbreak{}0} in {\ttfamily\footnotesize src/\allowbreak{}VERSION}.  The uploaded zip did not include a Git commit hash, so this manual records the archive content rather than a specific commit.  Generated tables were built from the C++ wrapper registrations, schema-generated RST files, ParticleFileWriter Python modules, and the suite directories.

\paragraph{What is included.} The manual focuses on {\ttfamily Solid\allowbreak{}Mechanics\_\allowbreak{}MPM}, MPM events, the particle mesh and particle-file format, ParticleFileWriter input generation, diagnostics, VisIt/ParaView output pathways, and the existing verification, validation, and example suite reports.  LLNL-specific external material-model reports are represented by a placeholder link until those private or out-of-tree reports are added to the documentation tree.

% AUTO-GENERATED FILE. DO NOT EDIT.
% Generated by tools/render_generated_tex.py.
{\renewcommand{\arraystretch}{1.15}
\small
\begin{longtable}{@{}p{0.27\linewidth}p{0.67\linewidth}@{}}
\caption{Code-derived content counts in this starter manual.}\label{tab:source-counts}\\
\toprule
\textbf{Topic} & \textbf{Notes}\\
\midrule
\endfirsthead
\toprule
\textbf{Topic} & \textbf{Notes}\\
\midrule
\endhead
\midrule
\multicolumn{2}{r}{Continued on next page}\\
\endfoot
\bottomrule
\endlastfoot
\textbf{Solver attributes} & 199\\
\textbf{MPM event types} & 19\\
\textbf{PFW parameters} & 205\\
\textbf{Particle file fields} & 44\\
\textbf{Geometry objects} & 42\\
\textbf{Geometry wrappers} & 24\\
\textbf{Post-processing scripts} & 73\\
\end{longtable}
}

